% Replace the original Overview/Methodology; merge the added analysis and validation.
% Copy packages, macros, and theorem declarations from risk_rf_revision.tex.
% Merge the supplied bibitems into the host bibliography.
\section{Overview}
\label{sec:overview}

The fusion problem is to regulate the influence of incomplete operational
reports before committing resources to adaptation. Workload records describe
task composition, operating telemetry describes execution conditions, and
service feedback describes the consequences of a deployed policy. Putting their
change scores in $[0,1]$ does not give them a common meaning. We therefore retain
source-specific diagnostics while calibrating each source to a shared,
subsequently observed service-loss target.

The proposed architecture has three connected stages. First, completed windows
provide source scores, validity masks, observation support, age, and sampling
variability. Second, forecasts issued before label arrival are combined using a
quality reference, their joint residual second moment, and a historical weight
anchor. Third, a decision layer distinguishes evidence of persistent loss from
evidence that a particular update can recover it. Candidate value must include
deployment delay, competing inference resources, and actual update cost.
Figure~\ref{fig:framework} shows this single, consistent pipeline.

\begin{figure}[htbp]
\centering
\small
\fbox{\begin{minipage}{0.93\linewidth}
\centering
\textbf{Completed operational windows}\\
Workload / operating telemetry / service loss; masks and timestamps\\[3pt]
$\Downarrow$\\[-1pt]
\textbf{Separate source diagnostics and common-target forecasts}\\
$d_{s,k}$; positive quality $q_{s,k}$; stored forecast $p_{s,k}$\\[3pt]
$\Downarrow$\\[-1pt]
\textbf{Quality-conditioned joint-risk fusion}\\
KL quality reference + full residual matrix + inertia + masked cap\\
$\Gamma_k,\widehat y_{k+1|k},D_k^p,c_k,r_k^{\rm pred}$\\[3pt]
$\Downarrow$\\[-1pt]
\textbf{Persistent-loss admission and candidate valuation}\\
Calibrated recovery / delay / uncertainty exposure / resource feasibility\\[3pt]
$\Downarrow$\\[-1pt]
\textbf{Executed inference and budgeted model updates}\\
Actual gradient steps, active costs, deployment times, and service outcomes\\[3pt]
$\curvearrowleft$ Stored forecasts receive delayed labels;
executed interventions calibrate recovery separately.
\end{minipage}}
\caption{Consistent DA-RF data flow. Joint prediction risk updates forecast
combination; executed interventions update a separate recovery model.
Neither a diagnostic change score nor agreement among sources establishes
that retraining is useful. All-missing evidence causes abstention.}
\label{fig:framework}
\end{figure}

Forecast combination, convex regularization, and dynamic source reliability
have established foundations~\cite{bates1969,bregman1967,boyd2004,diebold2019,
qmf2023,pdf2024,huang2025}. We do not claim that a KL projection, a quadratic
penalty, or correlated-error modeling is a new optimization primitive.
The methodological contribution is the specified interface between measurement
quality, delayed common-target risk, constrained influence, and recoverability-
aware decisions, together with tests that isolate their separate effects.

\section{Methodology}
\label{sec:methodology}

\subsection{Time, execution, and data provenance}
\label{sec:problem}

Time is divided into slots $t=1,\ldots,T$ and completed windows
$\mathcal T_k=\{(k-1)W+1,\ldots,kW\}$. An unfinished final window is logged
but cannot update fusion. Let $\mathcal F_t$ contain observations completed
before the slot-$t$ decision, plus current arrivals, capacities, and worker
status. A window-$k$ snapshot first becomes usable in slot $kW+1$.
All normalizers, thresholds, and initial models are fitted on a disjoint
historical prefix. Delayed labels update the forecasts originally saved at
their issuance times, consistent with delayed-feedback
learning~\cite{joulani2013}; this temporal discipline does not transfer an
external regret guarantee to our rule.

The coordinator chooses a training profile $i_t$ and an inference profile $j_t$.
The service policy chooses application actions within that inference allocation.
Its parameters change only when a completed worker is deployed. Let $s_t$ be
the completed fraction and $\ell_t$ the expired-or-rejected fraction of requests
present before service, with rejected arrivals included in both denominators.
Slots without requests are logged separately. A consistent external reward is
\begin{equation}
 r_t^{\rm env}=s_t-\alpha_\ell\ell_t-\alpha_q\bar q_t
                  -\alpha_C C_t^{\rm active},
 \label{eq:external-reward}
\end{equation}
where $\bar q_t$ is normalized post-service queue length. Lateness among completed
requests is a separate metric and is not denoted by $\ell_t$.
$C_t^{\rm active}$ charges each occupied training slot; a per-job total is
distributed over its occupied slots rather than charged in full at every slot.
No fusion score or recovery proxy appears in this evaluation reward.

All existing operational results used in this revision are synthetic. Alibaba
Cluster Trace v2018 is an optional input route, not the source of those tables.
A trace experiment requires actual file hashes, machine selection, time ranges,
aggregation rules, and a declaration of still-simulated quantities. A description
of an importer does not establish execution of a trace experiment.

\subsection{Source-specific evidence and a persistent common target}
\label{sec:evidence}

Each available source produces a diagnostic $d_{s,k}\in[0,1]$.
Workload divergence uses adjacent empirical task-regime distributions:
\begin{equation}
 d_{1,k}=\frac{\operatorname{JSD}(P_{k-1},P_k)}{\log 2}.
\end{equation}
Operating change uses prefix-normalized features:
\begin{equation}
 z_{t,l}=\frac{x_{t,l}-\mu_l^{\rm ref}}
                  {\max(\sigma_l^{\rm ref},\epsilon_\sigma)},
 \qquad
 d_{2,k}=\frac1{L_z}\sum_l\min\{1,|\bar z_{k,l}-\bar z_{k-1,l}|\}.
\end{equation}
Each mean uses documented valid rows. Service-change evidence compares
adjacent windows only on sufficiently supported common contexts
$h=(i,j,\text{deployed version})$. With
$\eta_{h,k}\propto\min(n_{h,k-1},n_{h,k})$ normalized over common contexts,
\begin{equation}
 d_{3,k}=\min\left\{1,
 \frac{[\bar r_{k-1}^{\rm match}-\bar r_k^{\rm match}]_+}
      {|\bar r_{k-1}^{\rm match}|+\epsilon_r}\right\}.
\end{equation}
The known direct training charge may be removed before this comparison;
resource-induced service loss is not thereby removed. Missing common support
makes this diagnostic unavailable, rather than zero.

An adjacent-window change can disappear while a damaged policy remains damaged.
DA-RF therefore predicts persistent loss rather than relying exclusively on the
positive adjacent-window difference. Let $J_k$ be cost-adjusted service utility
in window $k$ and let $J_k^{\rm ref}$ be a reference prediction fixed before
that window's outcomes arrive. The reference is fitted on the historical prefix
and conditions on declared workload, capacity, and inference-profile descriptors.
Define
\begin{equation}
 y_k=\clip\left(\frac{J_k^{\rm ref}-J_k}{B_J},0,1\right),\qquad B_J>0.
 \label{eq:persistent-target}
\end{equation}
Reference extrapolation, composition changes, and policy effects can confound
this target; it is an operational loss indicator, not an identified causal
drift probability. The reference fit and valid target contexts must be logged.
If the target cannot be computed, it is unavailable. The previous Risk RF
results used their original next-window degradation target; changing to
Eq.~\eqref{eq:persistent-target} requires a new experimental treatment.

Every source issues a forecast of the same subsequent target:
\begin{equation}
 p_{s,k}=\clip(\boldsymbol\theta_{s,k}^{\mathsf T}\boldsymbol\phi_{s,k},0,1),
 \qquad p_{s,k}\ \text{predicts }y_{k+1}.
 \label{eq:common-forecast}
\end{equation}
A minimal declared feature vector is
$\boldsymbol\phi_{s,k}=(1,d_{s,k},d_{s,k-1},y_k\ind\{y_k\text{ available}\},
\ind\{y_k\text{ available}\})^{\mathsf T}$; missing lag diagnostics receive
an explicit mask. The common service history is shared across branches and
baselines. Ridge-initialized exponentially weighted least squares updates
$\boldsymbol\theta_{s,k}$ only from earlier stored feature--target pairs whose
labels have arrived. Other predictors require the same time and feature budget.
Forecasts are frozen before their labels arrive and are never rescored after
refitting to those labels.

\subsection{Quality reference, delayed residual risk, and masked fusion}
\label{sec:reliability}

For available source $s$, observation quality is
\begin{equation}
 q_{s,k}=\max\left\{q_{\min},
 \frac{n_{s,k}}{n_{s,k}+n_0}
 \exp(-A_{s,k}/\tau_A)\frac1{\nu_{s,k}+\epsilon_\nu}\right\},
 \qquad q_{\min}>0.
 \label{eq:quality-only}
\end{equation}
Here $n_{s,k}$ is a valid support count, not an asserted independent effective
sample size; $A_{s,k}$ is the age of the most recent valid observation; and
$\nu_{s,k}$ is a moving-block-bootstrap variance of the diagnostic. Fewer than
two valid bootstrap replicates use a declared conservative variance fallback.
Availability thresholds precede the quality floor: the floor cannot make an
invalid source available. Measurement quality excludes prediction error.

When a complete stored forecast vector and its common label arrive, define
$\boldsymbol e_{k-1}=\boldsymbol p_{k-1}-y_k\boldsymbol1$ and update
\begin{equation}
 R_k=\beta R_{k-1}+(1-\beta)
             \boldsymbol e_{k-1}\boldsymbol e_{k-1}^{\mathsf T},
 \qquad R_0\succeq0,\quad0\le\beta<1.
 \label{eq:residual-update}
\end{equation}
This is an uncentered second moment: both common bias and correlated variation
contribute. The outer-product update preserves positive semidefiniteness.
Missing residuals are not replaced by zeros, and separate pairwise deletion is
not used because it can produce a non-PSD matrix. Incomplete vectors skip the
update and increment the risk-age counter. This complete-vector estimator can
be biased toward jointly observed windows; outage-conditioned error and update
counts must therefore be reported. Queue residual updates may additionally
require stable logged control contexts; this restriction and its rejection
rate must be declared. Estimated $R_k$ is not assumed equal to future population
risk.

Let $\mathcal A_k$ be the valid enabled sources, $N_k=|\mathcal A_k|$, and
$\bar a_k=\max(a_{\max},1/N_k)$. On this set,
$\pi_{s,k}=q_{s,k}/\sum_{v\in\mathcal A_k}q_{v,k}$.
The historical anchor $\boldsymbol a_k$ restricts and renormalizes previous
weights onto current available sources; if their retained mass is zero, it is
uniform. An anchor need not satisfy the current cap. The proposed fusion solves
\begin{align}
 \boldsymbol w_k&=\argmin_{\boldsymbol w\in\mathcal C_k}
 \left\{\tau D_{\rm KL}(\boldsymbol w\Vert\boldsymbol\pi_k)
       +\frac\lambda2\boldsymbol w^{\mathsf T}R_k\boldsymbol w
       +\frac\kappa2\norm{\boldsymbol w-\boldsymbol a_k}_2^2\right\},
 \label{eq:main-fusion}\\
 \mathcal C_k&=\{\boldsymbol w:\boldsymbol1^{\mathsf T}\boldsymbol w=1,
   0\le w_s\le\bar a_k,\ w_s=0\ (s\notin\mathcal A_k)\},
 \qquad \tau>0,\ \lambda,\kappa\ge0.
\end{align}
The available principal submatrix of $R_k$ remains PSD. With one source the
weight is one; no concentration protection is then claimed. With no sources
the module abstains, emits zero coverage and an invalid-forecast flag, and
launches no new training.

For three sources, an active-set Newton solver can enumerate the upper-cap
sets and solve the positive free-coordinate problem with an equality constraint.
For larger source sets, a convex solver returns the primal residual,
stationarity residual, tolerance, and iteration count. A reproducible setting
uses $10^{-8}$ primal feasibility and $10^{-7}$ stationarity tolerance, with
a declared iteration cap. A feasible capped-quality fallback is logged when
the solver fails; it is not silently treated as the optimizer.

The fusion state retains different interpretations:
\begin{align}
 \Gamma_k&=\sum_s w_{s,k}d_{s,k},&
 \widehat y_{k+1|k}&=\sum_s w_{s,k}p_{s,k},\\
 D_k^p&=\sum_s w_{s,k}(p_{s,k}-\widehat y_{k+1|k})^2,&
 c_k&=N_k/S_{\rm enabled},\qquad
 r_k^{\rm pred}=\boldsymbol w_k^{\mathsf T}R_k\boldsymbol w_k.
 \label{eq:fusion-state}
\end{align}
$\Gamma_k$ is diagnostic only. Disagreement compares forecasts of one target;
agreement alone does not establish accuracy, because all sources can share an
error. DA-RF includes joint risk and stale risk information in its exposure index:
\begin{equation}
 U_k=\min\{1,\omega_D4D_k^p+\omega_M(1-c_k)
           +\omega_R r_k^{\rm pred}+\omega_A A_k^{\rm risk}/A_{\max}\}.
 \label{eq:decision-uncertainty}
\end{equation}
The nonnegative coefficients are selected on historical validation data and
then frozen. This index is a decision feature, not a calibrated failure
probability or confidence interval. On source loss, forecasts and coverage
remain distinct; risk-estimator age cannot be reset without a genuine update.

\subsection{Original RF and matched fusion baselines}
\label{sec:baselines}

\textbf{Original RF is a baseline.} It multiplies measurement quality by
$q_{s,k}^E=\exp(-E_{s,k}/\tau_E)$, where $E_{s,k}$ is a delayed scalar
forecast MSE. Its normalized product is projected onto the capped simplex:
\begin{equation}
 \alpha_k^{\rm orig}=\argmin_{\boldsymbol\alpha\in\mathcal C_k}
 D_{\rm KL}(\boldsymbol\alpha\Vert
      \operatorname{normalize}(\boldsymbol q_k\odot\boldsymbol q_k^E)).
\end{equation}
It uses raw diagnostic disagreement, unlike the common-target disagreement
in Eq.~\eqref{eq:fusion-state}. Original RF has no joint residual matrix or
inertia term. Its equations and figures are not presented as the proposed rule.

\textbf{Diagonal-risk RF} replaces $R_k$ in Eq.~\eqref{eq:main-fusion} by
$\diag(R_k)$ while keeping every other input and hyperparameter unchanged.
This is the primary comparator for cross-source error structure. Removing the
entire risk penalty tests a different question. Equal and prefix-fixed fusion,
inverse marginal MSE, and declared online-expert rules receive the same stored
forecasts and delayed labels. Inverse MSE uses forecast residuals
$p_{s,k}-y_{k+1}$, not differences between heterogeneous diagnostics and labels.
An evidential baseline uses a declared mass construction and normalization,
and is not assigned fictitious linear weights. QMF and PDF require branch
predictions and are evaluated on compatible classification tasks; a core
adaptation is not a full reproduction of their original
systems~\cite{qmf2023,pdf2024}.

The fusion-only ablation holds the gate, target, forecast fitter, available
features, and recovery estimator fixed. The gate-only ablation holds forecasts
and weights fixed. This factorial separation prevents a new target or a more
permissive gate from being credited to the joint-risk matrix.

\subsection{Persistent-loss admission and recovery probes}
\label{sec:coordination}

At slot $t$, let $k(t)$ denote the latest completed window.
Define a persistent-loss statistic
\begin{equation}
 h_t=\max\{\widehat y_{k(t)+1|k(t)},y_{k(t)}^{\rm obs}\},
 \label{eq:loss-sentinel}
\end{equation}
where each term is included only if valid and within its declared age limit.
The observed-loss term is not a future label. With neither valid, adaptation
abstains. A hysteresis flag switches on at $h_t\ge\gamma_{\rm on}$ and off
at $h_t<\gamma_{\rm off}<\gamma_{\rm on}$. A large $\Gamma_k$ is not required:
a changed conditional relation can degrade service without changing workloads.

The no-update profile $i_0$ is always admitted when inference is resource feasible.
For substantial updates, the candidate set requires valid nonfallback evidence,
declared coverage, worker eligibility, and persistent-loss admission. Uncertainty
is charged in the candidate objective rather than used as an absolute veto that
could indefinitely suppress recovery after a loss.

When fresh observed service loss persists but update-value evidence is absent,
DA-RF may launch a small \emph{probe update}. Its maximum gradient steps,
cooldown, profile, active cost, and cumulative budget are fixed before evaluation.
A valid service-loss observation is mandatory; missing observations are not
interpreted as loss. The probe is an information-gathering exception, not a
certified beneficial intervention. Its cost enters the service endpoint, and an
otherwise identical observed-loss sentinel without fusion is an essential
baseline. Ablating the sentinel tests whether recovery depends on that mechanism
instead of fusion.

\subsection{Candidate-dependent uncertainty and measured recovery value}

Let $i\in\mathcal I$ index update profiles and $j\in\mathcal J$ inference
profiles. Residual capacity $B_{m,r,t}$ already excludes queued work and active
workers. For every node $m$ and resource $r$,
\begin{equation}
 D_t^{\rm tr}a^{\rm tr}_{i,m,r}
 +D_t^{\rm inf}a^{\rm inf}_{j,m,r}\le B_{m,r,t}.
 \label{eq:resource-feasibility}
\end{equation}
The base profile has zero update demand. Empty feasible sets invoke a declared
service admission/deferral policy rather than return an infeasible action.
Deployment delay $\delta_i$ and cost $C_t(i)$ represent actual executed or
explicitly simulated occupancy; wall-clock claims require hardware measurement.

Recovery is calibrated separately from forecast accuracy. From independent
calibration interventions, fork matched states into an update arm and a no-update
arm, preserve the same inference schedule and external random inputs, and record
the difference in gross service utility over the entire horizon, including
service losses before deployment and queue effects that persist afterward. Fit
$\widehat G_t(i,j)$ to the current fusion state, worker budget, profile, and
deployment delay using disjoint training and validation streams. Pending
deployments and the finite horizon are included. A bounded definition is
\begin{equation}
 G_t(i,j)=\E\!\left[
 \sum_{h=0}^{L_t}\gamma^{h}
  (u_{t+h}^{(i,j)}-u_{t+h}^{(i_0,j)})\mid\mathcal F_t\right],
 \label{eq:measured-recovery}
\end{equation}
where both rollout arms use the same declared inference protocol and the
update arm executes its actual resource occupancy. The policy parameters are
unchanged before deployment, but service can already change through resource
competition and queue dynamics; these predeployment differences are included.
The allowance horizon and discount are identical for service and costs.
In particular, $G_t(i_0,j)=\widehat G_t(i_0,j)=0$.
This estimand requires simulator forks or valid controlled interventions; it
cannot be recovered from arbitrary policy logs merely by naming an estimator
``causal.'' Recovery gains may be negative. Fixed profile constants are not
substituted for measured neural improvement. The prior finite-horizon recovery
proxy may be retained as a separate baseline but must be calibrated and audited.

Let $\widehat W_{j,t}$ predict the no-update gross utility over the same
finite horizon, including current service:
\begin{equation}
 W_{j,t}=\E\!\left[\sum_{h=0}^{L_t}\gamma^h
 u_{t+h}^{(i_0,j)}\mid\mathcal F_t\right].
\end{equation}
Let $C_t(i)$ be the discounted direct training charge over that horizon,
accounting for every active allocation slot. Training-induced completion,
deadline, and queue losses belong in $G_t(i,j)$, not merely in this direct
charge. Set $C_t(i_0)=0$. Let
$\epsilon_t(i,j)$ be a declared error allowance for
$\widehat W_{j,t}+\widehat G_t(i,j)$. It is not assumed to be a valid uniform
bound unless independently justified. With candidate intensity $Q_{i,t}$,
the decision score is
\begin{align}
 L_t(i,j)={}&\widehat W_{j,t}+\widehat G_t(i,j)-\epsilon_t(i,j)
             -\lambda_C C_t(i)\nonumber\\
            &-\omega_Q Q_{i,t}^2-\lambda_U U_t Q_{i,t},
 \qquad Q_{i_0,t}=0.
 \label{eq:candidate-score}
\end{align}
Unlike a common subtraction $-\lambda_U U_t$, the last term changes candidate
comparisons. For two candidates, its contribution is
$-\lambda_U U_t(Q_{i,t}-Q_{i',t})$; all else equal, greater uncertainty favors
lower resource exposure. It has no selection effect if every intensity is equal,
which must be recorded as an inactive test condition.

The no-update upper reference is
\begin{equation}
 B_t^0=\max_{j:(i_0,j)\ \rm feasible}
       \{\widehat W_{j,t}+\epsilon_t(i_0,j)\}.
\end{equation}
For exploitation updates, admit only candidates with
$L_t(i,j)>B_t^0+m_t$, $m_t\ge0$, and maximize Eq.~\eqref{eq:candidate-score}
over the admitted feasible pairs. The separately budgeted probe exception is
logged. Otherwise select a feasible no-update inference profile. Exact
enumeration is inexpensive for five training and six inference profiles and
is optimal for the declared score only. Larger profile libraries may use
alternating optimization with a measured enumeration gap; coordinate convergence
does not establish global optimality.

\begin{proposition}[Conditional net-value certificate]
\label{prop:recovery-certificate}
Suppose every candidate's true finite-horizon gross value
$W_{j,t}+G_t(i,j)$ differs from its prediction by at most
$\epsilon_t(i,j)$, and costs are correctly measured. Then an exploitation
candidate satisfying $L_t(i,j)>B_t^0+m_t$ has true net value exceeding that
of every feasible no-update candidate by more than $m_t$. This statement does
not apply to exploratory probes or to uncertified empirical error allowances.
\end{proposition}
\begin{proof}
The true net value is at least
$\widehat W_{j,t}+\widehat G_t(i,j)-\epsilon_t(i,j)-\lambda_C C_t(i)$,
which is at least $L_t(i,j)$ because the exposure penalties are nonnegative.
Every no-update true value is at most $B_t^0$. The admission inequality gives
the result. It concerns the specified finite-horizon comparison and does not
establish long-run policy regret under changing unmodeled conditions.
\end{proof}

The certificate identifies what must be verified rather than guaranteeing that
such a candidate exists. Recoverability requires an actual update profile that
improves the deployed policy sufficiently to cover delay and cost. Persistent
loss and probe availability prevent the particular logical failure of demanding
a new adjacent-window change before investigating a still-damaged policy;
they do not guarantee optimizer convergence or policy recovery.

\subsection{Execution order and computational audit}

At each slot: (1) ingest current capacities and the last valid fusion snapshot;
(2) update the persistent-loss flag using completed feedback;
(3) form resource-feasible no-update, exploitation, and budgeted-probe candidates;
(4) compute candidate values once and enumerate the admitted pairs;
(5) execute service, charge active costs, and record worker launch/completion;
(6) deploy only completed workers; and (7) on a completed window, issue and
store forecasts, update earlier stored forecasts whose labels have arrived,
compute quality and joint risk, and solve fusion for the next slot.
When issuing a forecast and fitting from an older arrived target occur at the
same boundary, the declared sequence uses only targets already available before
the new forecast is issued. Forecast identifiers make this order auditable.

The timing report separates source extraction and bootstrap, predictor fitting,
residual updates, convex fusion, candidate-value estimation, and pair search.
An outer-product update costs $O(S^2)$ and stores $O(S^2)$ moments. For three
sources, solving the convex fusion is small, but bootstrap and neural updates
can dominate total cost. A submillisecond pair-search timing is not reported as
end-to-end fusion or training latency.

\section{Theoretical analysis}
\label{sec:theory}
In the next results, $R$ is a generic PSD input matrix. When interpreting
$\boldsymbol w^{\mathsf T}R\boldsymbol w$ as population prediction error,
$R$ must be the true second moment for the relevant test distribution. An
estimated moving-window matrix needs a separate estimation-error argument.
The symbol $b$ denotes a fixed cap in the stability analysis, while $b_0$
denotes the target half-width in the threshold-decision construction.

% Independently checked theory fragment; integrate notation into the main manuscript.
% These are theoretical statements, not measured experimental results.

\subsection{Stability and the incremental value of joint residual risk}
\label{sec:joint-risk-theory}

The joint-risk model requires comparable predictions: every available source
produces a causally calibrated estimate $p_{s,k}$ of the same scalar target
$z_k$. With $e_{s,k}=p_{s,k}-z_k$, the matrix
$R_k=\mathbb E[\boldsymbol e_k\boldsymbol e_k^{\mathsf T}]$ is the residual
second moment, including both bias and covariance. Consequently,
\begin{equation}
 \mathbb E\!\left[(\boldsymbol w^{\mathsf T}\boldsymbol p_k-z_k)^2\right]
 =\boldsymbol w^{\mathsf T}R_k\boldsymbol w.
 \label{eq:joint-risk-identity}
\end{equation}
This identity does not apply to an uncalibrated average of workload divergence,
resource variation, and service loss. These heterogeneous quantities must first
be mapped to the shared target using only previously completed outcomes.

Fix an availability mask with $n$ available sources and a cap
$b\in[1/n,1]$. Let
\begin{equation}
 \mathcal C_b=\left\{\boldsymbol w\in\mathbb R^n:
 \boldsymbol 1^{\mathsf T}\boldsymbol w=1,\quad 0\le w_s\le b\right\},
 \qquad
 \pi_s(\boldsymbol q)=\frac{q_s}{\sum_r q_r},\quad q_s>0.
\end{equation}
The regularized full-risk fusion rule is
\begin{equation}
 \boldsymbol w(\boldsymbol q,R,\boldsymbol a)=
 \arg\min_{\boldsymbol w\in\mathcal C_b}
 \left\{
 \tau D_{\mathrm{KL}}(\boldsymbol w\Vert\boldsymbol\pi(\boldsymbol q))
 +\frac{\lambda}{2}\boldsymbol w^{\mathsf T}R\boldsymbol w
 +\frac{\kappa}{2}\|\boldsymbol w-\boldsymbol a\|_2^2
 \right\},
 \label{eq:checked-fusion-objective}
\end{equation}
where $R\succeq0$, $\tau>0$, and $\lambda,\kappa\ge0$.
The anchor $\boldsymbol a$ is fixed while solving the current window; it can be
the previous weights restricted and renormalized to the available sources.

\begin{proposition}[Existence, uniqueness, and sensitivity]
\label{prop:fusion-sensitivity}
The solution of Eq.~\eqref{eq:checked-fusion-objective} exists and is unique.
For two parameter triples with the same mask, cap, and regularization
coefficients, define
$\delta_q=\min_{c\in\mathbb R}
 \|\log\boldsymbol q-\log\widetilde{\boldsymbol q}-c\boldsymbol1\|_2$.
Their solutions obey
\begin{equation}
 \|\boldsymbol w-\widetilde{\boldsymbol w}\|_2
 \le
 \frac{
  \tau\delta_q+
  \lambda\sqrt b\,\|R-\widetilde R\|_{\mathrm{op}}+
  \kappa\|\boldsymbol a-\widetilde{\boldsymbol a}\|_2
 }{\tau/b+\kappa}.
 \label{eq:three-way-sensitivity}
\end{equation}
In particular, replacing $\delta_q$ by
$\|\log\boldsymbol q-\log\widetilde{\boldsymbol q}\|_2$ and $\sqrt b$ by
$1$ gives a weaker valid bound. Equation~\eqref{eq:three-way-sensitivity}
does not compare windows whose availability masks or caps differ.
\end{proposition}

\begin{proof}
The objective is continuous on the nonempty compact set $\mathcal C_b$, with
$0\log0=0$. On its positive coordinates its Hessian is
\begin{equation}
 H=\tau\operatorname{diag}(1/w_s)+\lambda R+\kappa I
 \succeq(\tau/b+\kappa)I.
\end{equation}
The same strong-convexity inequality extends to the boundary by continuity,
proving uniqueness. The optimizer has positive weights: for $b>1/n$, mixing
a point having a zero coordinate with the uniform vector has an entropy
directional derivative of $-\infty$, whereas the other terms have finite
directional derivatives. For $b=1/n$, the feasible set consists only of the
strictly positive uniform vector.

Write $m=\tau/b+\kappa$ and
$\boldsymbol\delta=\boldsymbol w-\widetilde{\boldsymbol w}$.
The two variational inequalities and strong convexity give
\begin{align}
 m\|\boldsymbol\delta\|_2^2
 &\le
 \left[
  \tau(\log\boldsymbol\pi-\log\widetilde{\boldsymbol\pi})
  +\lambda(\widetilde R-R)\widetilde{\boldsymbol w}
  +\kappa(\boldsymbol a-\widetilde{\boldsymbol a})
 \right]^{\mathsf T}\boldsymbol\delta.
\end{align}
Since $\boldsymbol1^{\mathsf T}\boldsymbol\delta=0$, the log-normalization
constants and any constant multiple of $\boldsymbol1$ vanish. Moreover,
$\|\widetilde{\boldsymbol w}\|_2^2\le b\sum_s\widetilde w_s=b$.
Cauchy--Schwarz completes the proof.
\end{proof}

The bound separates three perturbations. Quality changes enter through
relative log-quality; a common multiplicative rescaling of all $q_s$ has no
effect. Residual estimation error enters through an operator norm and can
change weights even when all marginal residual variances remain unchanged.
An anchor perturbation is attenuated by the factor
$\kappa/(\tau/b+\kappa)$. Increasing $\kappa$ improves temporal stability but
can slow redistribution after a persistent regime change. If quality is a
product of support, freshness, and variability factors, their log changes
add in the quality term; all factors require a strictly positive floor.

When no cap is active, a more informative local expression follows by
differentiating the first-order conditions. Define
\begin{equation}
 Q=H^{-1}-H^{-1}\boldsymbol1
 (\boldsymbol1^{\mathsf T}H^{-1}\boldsymbol1)^{-1}
 \boldsymbol1^{\mathsf T}H^{-1}.
\end{equation}
Then
\begin{equation}
 d\boldsymbol w
 =Q\left(\tau\,d\log\boldsymbol q
         -\lambda\,(dR)\boldsymbol w
         +\kappa\,d\boldsymbol a\right).
 \label{eq:local-sensitivity}
\end{equation}
For a fixed active set, the same expression applies to its free coordinates;
coordinates held at a bound have zero differential. At an active-set change,
Eq.~\eqref{eq:three-way-sensitivity} remains applicable, whereas the local
Jacobian need not be differentiable.

\begin{proposition}[When joint structure changes the fusion rule]
\label{prop:joint-structure-criterion}
Let $R_D=\operatorname{diag}(R)$ and $O=R-R_D$. Keep quality, anchor, cap,
and regularization coefficients fixed, with $\lambda>0$. Suppose that the
diagonal-risk solution $\boldsymbol w_D$ lies strictly below every cap. Then
it is also the full-risk solution if and only if
\begin{equation}
 \left(I-\frac{\boldsymbol1\boldsymbol1^{\mathsf T}}{n}\right)
 O\boldsymbol w_D=\boldsymbol0.
 \label{eq:joint-structure-condition}
\end{equation}
\end{proposition}

\begin{proof}
All solution coordinates are positive. The diagonal objective is stationary
on the simplex at $\boldsymbol w_D$. The full objective differs in gradient
by $\lambda O\boldsymbol w_D$. Thus, it is stationary at the same point
exactly when this difference is a constant multiple of $\boldsymbol1$.
Strict convexity makes stationarity sufficient and the solution unique.
\end{proof}

This criterion rules out a universal improvement claim: an equicorrelated
matrix with identical marginal risks, uniform quality, and a uniform anchor
leaves the uniform weights unchanged.
The useful case is unequal redundancy, where some sources share errors and
another provides complementary information. The following construction
isolates this effect without changing any marginal mean-square error.

\begin{proposition}[Complementarity at matched marginal error]
\label{prop:matched-marginal-gain}
Consider three unbiased sources with
\begin{equation}
 R(\rho)=\sigma^2
 \begin{bmatrix}1&\rho&0\\\rho&1&0\\0&0&1\end{bmatrix},
 \qquad \sigma^2>0,\quad 0<\rho<1,
 \label{eq:matched-marginal-matrix}
\end{equation}
uniform quality and anchor
$\boldsymbol u=(1/3,1/3,1/3)^{\mathsf T}$, and a cap $b\ge1/2$.
Every source has marginal mean-square error $\sigma^2$.
The diagonal-risk rule chooses $\boldsymbol u$.
Without quality or anchor regularization, the full-risk minimizer is
\begin{equation}
 \boldsymbol w^*(\rho)=
 \frac{(1,1,1+\rho)^{\mathsf T}}{3+\rho},
 \qquad
 \mathcal R(\boldsymbol u)-\mathcal R(\boldsymbol w^*)
 =\frac{2\sigma^2\rho^2}{9(3+\rho)}>0,
 \label{eq:matched-marginal-gap}
\end{equation}
where $\mathcal R(\boldsymbol w)=\boldsymbol w^{\mathsf T}R(\rho)\boldsymbol w$.
With any finite $\tau>0$, $\kappa\ge0$, and $\lambda>0$, the full-risk
solution still has strictly smaller true residual risk than the
diagonal-risk rule.
\end{proposition}

\begin{proof}
The matrix is positive definite for $0<\rho<1$. Symmetry gives
$\boldsymbol w=((1-y)/2,(1-y)/2,y)^{\mathsf T}$, and
\begin{equation}
 \mathcal R(y)=\sigma^2\left\{
 \frac{1+\rho}{2}(1-y)^2+y^2\right\}.
\end{equation}
Its minimizer is $y^*=(1+\rho)/(3+\rho)$, yielding
$\mathcal R(\boldsymbol w^*)=\sigma^2(1+\rho)/(3+\rho)$ and
$\mathcal R(\boldsymbol u)=\sigma^2(3+2\rho)/9$.
For the regularized rule, the unique $y_\rho$ solves
\begin{equation}
 \frac{\lambda\sigma^2}{2}
 \big[(3+\rho)y_\rho-(1+\rho)\big]
 +\tau\log\frac{2y_\rho}{1-y_\rho}
 +\frac{3\kappa}{2}\left(y_\rho-\frac13\right)=0.
 \label{eq:regularized-three-source-root}
\end{equation}
The left-hand side is strictly increasing, negative at $1/3$, and positive
at $y^*$. Therefore $1/3<y_\rho<y^*<1/2$; the cap is inactive.
Risk decreases strictly over this interval. More explicitly, with
$\delta=y_\rho-1/3$,
\begin{equation}
 \mathcal R(\boldsymbol u)-\mathcal R(y_\rho)
 =\sigma^2\left[
 \frac{2\rho}{3}\delta-\frac{3+\rho}{2}\delta^2\right]>0.
 \label{eq:regularized-risk-gain}
\end{equation}
\end{proof}

For this construction, define
\begin{equation}
 L_\rho=\frac{\lambda\sigma^2}{2}(3+\rho)
 +\tau\left(\frac1{y_\rho}+\frac1{1-y_\rho}\right)
 +\frac{3\kappa}{2}>0.
\end{equation}
Implicit differentiation gives
\begin{align}
 \frac{d y_\rho}{d\rho}
 &=\frac{\lambda\sigma^2(1-y_\rho)}{2L_\rho}>0,\\
 \frac{d y_\rho}{d\tau}
 &=-\frac{\log[2y_\rho/(1-y_\rho)]}{L_\rho}<0,\\
 \frac{d y_\rho}{d\kappa}
 &=-\frac{3(y_\rho-1/3)}{2L_\rho}<0.
 \label{eq:correlation-regularization-sensitivity}
\end{align}
Thus, stronger paired-error correlation transfers weight toward the
independent source, while the quality reference and historical anchor
attenuate that transfer. Equation~\eqref{eq:matched-marginal-gap} is a
theoretical construction, not an empirical result.

A quantitative sufficient condition also accounts for matrix estimation.
Let $h(\boldsymbol w)=\tau D_{\mathrm{KL}}(\boldsymbol w\Vert\boldsymbol u)
+\kappa\|\boldsymbol w-\boldsymbol u\|_2^2/2$ and let $\widehat R\succeq0$
satisfy $\|\widehat R-R\|_{\mathrm{op}}\le\epsilon_R$.
For its regularized solution $\widehat{\boldsymbol w}$,
\begin{equation}
 \mathcal R(\boldsymbol u)-\mathcal R(\widehat{\boldsymbol w})
 \ge
 \frac{2\sigma^2\rho^2}{9(3+\rho)}
 -\frac{2h(\boldsymbol w^*)}{\lambda}-2b\epsilon_R.
 \label{eq:estimable-risk-gain}
\end{equation}
Indeed, compare the estimated objective at
$\widehat{\boldsymbol w}$ and $\boldsymbol w^*$, use $h\ge0$, and bound
each quadratic estimation error by $b\epsilon_R$.
Here
\begin{align}
 D_{\mathrm{KL}}(\boldsymbol w^*\Vert\boldsymbol u)
 &=\frac{2}{3+\rho}\log\frac{3}{3+\rho}
 +\frac{1+\rho}{3+\rho}\log\frac{3(1+\rho)}{3+\rho},\\
 \|\boldsymbol w^*-\boldsymbol u\|_2^2
 &=\frac{2\rho^2}{3(3+\rho)^2}.
\end{align}
The lower bound is informative only when its right-hand side is positive.
It identifies regularization bias and joint-moment estimation error as
distinct reasons why the complementary-source advantage may not be observed.

These fusion-level statements do not establish that retraining recovers a
service policy or that weight changes cross a scheduling threshold. Those
claims require executed-update experiments with matched training budgets,
deployment timing, and drift events. Both weight/action sensitivity and
observed recovery must be reported; unchanged actions in an ablation must
be retained as a negative result.

\begin{proposition}[An exact fusion-to-decision bridge]
\label{prop:threshold-decision-bridge}
Let $z$ be uniform on $[\theta-b_0,\theta+b_0]$ and independent of the
entire error vector $\boldsymbol e$. Suppose
$|e_s|\le\sigma<b_0$ almost surely,
$\theta-b_0-\sigma\ge0$, and $\theta+b_0+\sigma\le1$.
Each source reports $p_s=z+e_s$.
Let $\boldsymbol w$ be a simplex weight vector selected independently of
the test pair $(z,\boldsymbol e)$ before the test batch. Define the executed
threshold action and utility by
\begin{equation}
 a(\boldsymbol w)=\mathbb I\{\boldsymbol w^{\mathsf T}\boldsymbol p\ge\theta\},
 \qquad J(z,a)=a(z-\theta),
 \qquad a^*(z)=\mathbb I\{z\ge\theta\}.
\end{equation}
With $R=\mathbb E[\boldsymbol e\boldsymbol e^{\mathsf T}]$, the expected
executed-decision regret satisfies the exact identity
\begin{equation}
 \mathbb E[J(z,a^*)-J(z,a(\boldsymbol w))\mid\boldsymbol w]
 =\frac{\boldsymbol w^{\mathsf T}R\boldsymbol w}{4b_0}.
 \label{eq:exact-decision-regret}
\end{equation}
No zero-mean error assumption is required.
\end{proposition}

\begin{proof}
Condition on $\boldsymbol w$ and $\boldsymbol e$, and write
$x=z-\theta$ and $\eta=\boldsymbol w^{\mathsf T}\boldsymbol e$.
The simplex constraint gives $|\eta|\le\sigma<b_0$.
For $\eta>0$, a suboptimal action occurs only for $-\eta\le x<0$, with
regret $-x$. For $\eta<0$, it occurs only for $0\le x<-\eta$, with regret
$x$. Thus, in either case,
\begin{equation}
 \mathbb E[J(z,a^*)-J(z,a(\boldsymbol w))\mid\boldsymbol w,\boldsymbol e]
 =\frac{\eta^2}{4b_0}.
\end{equation}
The same identity holds for $\eta=0$. Averaging over the independent test
errors yields Eq.~\eqref{eq:exact-decision-regret}.
\end{proof}

The matched-marginal matrix in Eq.~\eqref{eq:matched-marginal-matrix} can
be realized exactly with bounded errors. Let $A$ and $C$ be independent
symmetric Rademacher variables and let $S$ be independent of them, with
$\Pr(S=1)=(1+\rho)/2$. Set
\begin{equation}
 \boldsymbol e=\sigma(A,AS,C)^{\mathsf T}.
\end{equation}
All source marginals are identical, while
$\mathbb E[e_1e_2]=\sigma^2\rho$ and the other cross moments vanish.
Consequently, the strict regularized full-risk improvement in
Proposition~\ref{prop:matched-marginal-gain} translates into strictly
higher expected utility for this executed threshold decision. In the
unregularized case the utility gain is
\begin{equation}
 \mathbb E[J(z,a(\boldsymbol w^*))-J(z,a(\boldsymbol u))]
 =\frac{\sigma^2\rho^2}{18b_0(3+\rho)}.
\end{equation}
A concrete bounded setting is $\theta=0.5$, $b_0=0.25$, and
$\sigma=0.15$, for which every source prediction lies in $[0.1,0.9]$.
This mechanism benchmark uses known moments and an independent test batch.
It establishes an exact fusion-to-decision link in its stated setting;
it does not establish policy-update recovery in the queueing service.


\subsection{Propagation to admission and executed profile selection}
\label{sec:downstream-sensitivity}

Consider two parameter settings with identical source forecasts,
diagnostics, masks, cap, resource-feasible candidates, worker state, and
regularization coefficients. Let $B_w$ be the right-hand side of
Eq.~\eqref{eq:three-way-sensitivity}, and write
$\Delta R=R-\widetilde R$. Since simplex weight differences sum to zero,
\begin{align}
 |\Gamma-\widetilde\Gamma|
 &\le\min_{a\in\R}\norm{\boldsymbol d-a\boldsymbol1}_2 B_w,\\
 |\widehat y-\widetilde{\widehat y}|
 &\le\min_{a\in\R}\norm{\boldsymbol p-a\boldsymbol1}_2 B_w,\\
 |D^p-\widetilde D^p|&\le\tfrac12\norm{\boldsymbol w-
                      \widetilde{\boldsymbol w}}_1
                      \le\tfrac{\sqrt n}{2}B_w,\\
 |r^{\rm pred}-\widetilde r^{\rm pred}|
 &\le2\sqrt b\norm{R}_{\rm op}B_w+b\norm{\Delta R}_{\rm op}.
 \label{eq:representation-sensitivity}
\end{align}
For the disagreement bound, along the line between two simplex vectors the
directional derivative is
$\sum_s\Delta w_s(p_s-\bar p)^2$. The squared terms lie in $[0,1]$ and
$\sum_s\Delta w_s=0$, giving the half-$\ell_1$ bound after integration.
The risk bound follows by expanding the two quadratic forms and using
$\norm{\boldsymbol w}_2,\norm{\widetilde{\boldsymbol w}}_2\le\sqrt b$.
With coverage and risk age held fixed, clipping is nonexpansive, so
\begin{equation}
 |U-\widetilde U|\le
 2\omega_D\sqrt n B_w+
 \omega_R\{2\sqrt b\norm{R}_{\rm op}B_w+b\norm{\Delta R}_{\rm op}\}.
 \label{eq:uncertainty-sensitivity}
\end{equation}
These inequalities hold for fixed forecasts. Refitting predictors introduces
additional forecast perturbations and must not be silently attributed to a
weight-only perturbation.

For fixed current inference and cost estimates, candidate-score changes obey
\begin{equation}
 \Delta L_t(i,j)=\Delta\widehat G_t(i,j)
                -\lambda_UQ_{i,t}\Delta U_t-\Delta\epsilon_t(i,j).
 \label{eq:score-propagation}
\end{equation}
Suppose the recovery estimator is Lipschitz on the declared operating domain
and the resulting absolute score perturbation is at most $B_L$ for every
candidate. If the winning pair has a score gap greater than $2B_L$, its identity
cannot change. The same margin argument applies to an exploitation-admission
comparison with $B_t^0$. A gate or active-set boundary requires explicit
analysis rather than a differentiable Jacobian argument.

This result explains why a valid component may not change an action in a
particular experiment. Conversely, a meaningful component test must include
states where its perturbation changes source influence and where the candidate
margin is small enough for that influence to affect admission or intensity.
Both changed and unchanged actions are reported. We do not choose test states
after observing a favorable service outcome or infer a service gain from a
weight change alone.

\section{Known-moment mechanism verification}
\label{sec:known-moment-verification}

The following small experiment was executed for this revision using the supplied
\texttt{verify\_joint\_error.py}. It verifies the theoretical construction,
not the DA-RF online adaptation policy. Each of twenty independent seeds
contains $100{,}000$ target--error pairs per correlation condition. The target
is uniform on $[0.25,0.75]$; $\theta=0.5$, $b_0=0.25$, and $\sigma=0.15$.
The bounded Rademacher construction gives exactly the same marginal squared
error $\sigma^2$ for all three sources on every trial. Test targets are
independent of the complete error vector.

The diagonal and full rules share uniform quality and anchor, cap $0.7$,
$\tau=0.001$, $\lambda=1$, and $\kappa=0.001$. These coefficients are declared
mechanism settings, not selected on queue-service evaluation data. The true
matrix $R(\rho)$ is supplied to the rules, and the regularized full optimum is
computed by bisection of Eq.~\eqref{eq:regularized-three-source-root}.
An unregularized known-moment oracle is saved separately. Every method executes
its own threshold action on identical test pairs; utility and regret are computed
from the actual selected action. There is no delayed moment estimation, model
training, resource queue, or probe intervention in this verification.

Table~\ref{tab:mechanism-verification} reports mean executed-decision regret
and the paired full-minus-diagonal utility difference. The population identity
in Eq.~\eqref{eq:exact-decision-regret} uses $4b_0=1$, so population regret
equals population fused MSE in this particular setting. Zero correlation is
the negative control: the two rules coincide. Correlation among the first two
sources produces a conditional decision benefit without changing their marginal
errors. This does not establish a benefit on arbitrary error distributions.

% The measured table is inserted by the source assembler from the saved CSV.

\begin{table}[htbp]
\centering\small
\caption{Newly executed known-moment threshold-decision verification. Entries are mean $\pm$ sample SD across twenty seeds. This is a controlled mechanism example, not an edge-service experiment.}
\label{tab:mechanism-verification}
\resizebox{\linewidth}{!}{%
\begin{tabular}{rrrr}
\toprule
$\rho$ & Diagonal regret & Full regularized regret & Paired utility gain\\
\midrule
0.00 & $0.0074820\pm0.0000795$ & $0.0074820\pm0.0000795$ & $0.0000000\pm0.0000000$\\
0.25 & $0.0087267\pm0.0000915$ & $0.0086302\pm0.0000941$ & $0.0000965\pm0.0000191$\\
0.50 & $0.0099751\pm0.0000903$ & $0.0096191\pm0.0000926$ & $0.0003560\pm0.0000293$\\
0.75 & $0.0112259\pm0.0000832$ & $0.0104809\pm0.0000851$ & $0.0007450\pm0.0000311$\\
0.90 & $0.0119784\pm0.0000848$ & $0.0109498\pm0.0000872$ & $0.0010286\pm0.0000229$\\
\bottomrule
\end{tabular}}
\end{table}

A separate numerical audit solves one hundred paired convex problems with
perturbed positive quality, PSD residual matrices, and historical anchors while
holding the mask, cap, and regularization fixed. Every observed weight change
is below the bound in Eq.~\eqref{eq:three-way-sensitivity}; the largest observed
change-to-bound ratio is $0.63412$ (rounded upward). The numerical tolerance is
$2\times10^{-6}$. This finite audit supports implementation consistency and
does not replace the proof. The saved CSV records the predeclared settings,
weights, population risks, empirical MSE, executed utility, and regret; the JSON
records the random seed protocol and script hash.

\section{Existing evidence and validation required for the extension}
\label{sec:validation}

% Integration fragment: no document class or package declarations.
% Requires the main document's usual amsmath support.
% Existing evidence and proposed experiments are intentionally separated.
% All protocols below are prespecified designs; no new results are reported.

\subsection{Evidence boundaries and prespecified validation}
\label{sec:validation-boundaries}

\paragraph{Evaluated implementation versus proposed extension.}
We use \emph{Risk RF} exclusively for the implementation that generated the
reported predictions, admissions, and executed outcomes. Its weights incorporate
measurement quality, a residual second-moment penalty, and temporal inertia;
its admission rule uses common-target forecast disagreement and source coverage.
Any subsequently introduced decision-aware extension is denoted
\emph{DA-RF} and is an unevaluated method proposal in the present manuscript.
Changing an admission rule, candidate-value estimator, reward definition,
feedback filter, or deployment mechanism creates a different experimental
treatment. Existing Risk RF results therefore cannot establish the performance
of DA-RF. The protocols below specify the evidence needed to evaluate that
extension; they contain no newly executed experiments or numerical outcomes.

\paragraph{What the existing results establish.}
The observer replay supports specific evidence-level effects. Under operating
noise, Risk RF admits $10/10$ scheduled event trials compared with $8/10$ for
Original RF, while both have a zero nominal raw-alarm rate. This gain has a
boundary: under service-feedback outage, admitted coverage falls from $10/10$
for Original RF to $8/10$ for Risk RF. Under workload conflict, both admit
$10/10$ events, but nominal admission increases from $5.96\%$ to $9.47\%$.
Removing variability weighting increases the noise-induced raw-alarm rate
from zero to $8.42\%$ for Original RF and $7.02\%$ for Risk RF. These results
support conditional effects on weighting and admission; they do not establish
uniform superiority or a downstream service gain.

The joint-shift queue experiment separates allocation from updating. With
adaptive inference and no retraining, completion is $80.40\%$ for DQN and
$81.19\%$ for PPO, versus $59.59\%$ and $59.83\%$ with a fixed inference
profile. The respective gains of $20.81$ and $21.36$ percentage points are
attributable to the allocation comparison. Risk RF completion is $80.40\%$
and $81.17\%$, and its returns of $1799.10$ and $1799.14$ remain below the
adaptive no-retraining returns of $1799.95$ and $1801.46$. These observations
do not demonstrate an incremental benefit from RF-triggered updating.

The hidden-relation diagnostic provides an additional negative result.
DQN Risk RF obtains post-shift accuracy $3.43\%$ and final-512 accuracy
$3.36\%$, matching the no-retraining control. Observed-loss admission obtains
$55.76\%$ and $93.52\%$, respectively. Thus, $93.52\%$ is an observed-loss
baseline result, not a recovery result of the proposed fusion trigger.
PPO Risk RF obtains $0.17\%$ post-shift accuracy and $0.20\%$ final-512
accuracy, compared with $17.23\%$ and $40.74\%$ under observed-loss admission.
The diagnostic shows that an updating intervention can help under the specified
DQN configuration, while the evaluated fusion admission does not deliver that
recovery. A common maximum budget does not imply equal realized expenditure.

In the independent delayed-label classification adaptation, Risk RF accuracy is
$(96.26\pm2.08)\%$, compared with $(94.80\pm2.40)\%$ for equal-logit fusion
and $(94.45\pm2.44)\%$ for PDF. The mean differences of $1.46$ and $1.81$
percentage points require paired uncertainty estimates before stronger
comparative claims are made. Risk RF NLL/ECE are $0.5424/0.3669$, worse than
equal-logit fusion's $0.5267/0.3366$ and PDF's $0.4466/0.2715$. Accuracy gain
therefore does not establish improved predictive calibration.

\paragraph{Provenance and statistical scope.}
The reported numerical studies use controlled synthetic inputs. An optional
trace-import interface, or a description of Alibaba data fields, is not evidence
that a reported table used trace-derived inputs. A future trace experiment must
identify the actual files, hashes, machines, time interval, aggregation and
normalization rules, and the quantities that remain simulated. No current result
is relabeled as trace-based. Five paired seeds and ten scheduled event trials
provide a limited descriptive assessment. Resampling those seeds does not add
independent experimental units. New studies will save per-seed paired
differences, expose the event denominator, and separate scenario-specific
effects from pooled averages.

\subsubsection{Isolating the incremental contribution of cross-source error}
\label{sec:validation-cross-error}

The primary forecast experiment will compare the full residual matrix $R_k$
against $\operatorname{diag}(R_k)$ with identical source forecasts, delayed
targets, quality references, caps, inertia, initialization, and historical
hyperparameter selection. Equal weighting, inverse marginal MSE, and the
declared online-expert aggregation controls will receive the same observations.
This comparison isolates cross-error information rather than a change in source
precision, forecast fitting, or admission thresholds.

A controlled generator will keep each source's marginal squared error identical
while changing its relationship to another source. Let $y_{k+1}=1/2$ and
$0<\sigma\leq1/4$, and generate stored forecasts
\begin{align}
 p_{1,k}&=y_{k+1}+\sigma z_k,&
 p_{2,k}&=y_{k+1}+\sigma b_kz_k,&
 p_{3,k}&=y_{k+1}+\sigma v_k,
 \label{eq:validation-error-generator}
\end{align}
where $z_k$ and $v_k$ are independent equiprobable signs, $b_k$ is independent
of both, and
$\Pr(b_k=1)=(1+\rho)/2$. The parameter $\rho$ takes prespecified positive,
zero, and negative values, including $0.8$, $0$, and $-0.8$.
Every source has squared error exactly $\sigma^2$ on every trial; no clipping
is needed. The expected second-moment matrix is
\begin{equation}
 R(\rho)=\sigma^2
 \begin{pmatrix}1&\rho&0\\ \rho&1&0\\0&0&1\end{pmatrix}.
 \label{eq:validation-error-matrix}
\end{equation}
Positive $\rho$ represents redundant failure and negative $\rho$ represents
error cancellation. The source-quality reference is fixed and uniform in this
isolation experiment. Targets are withheld until the declared delay expires;
the fusion implementation receives neither future signs nor generator parameters.
The generator is not a fitted model of a real operational source.

As an analytical reference, with no entropy or inertia penalty and a nonbinding
cap, the minimum-risk simplex weights are
\begin{equation}
 w_1^*=w_2^*=\frac{1}{3+\rho},\qquad
 w_3^*=\frac{1+\rho}{3+\rho}.
 \label{eq:validation-oracle-weights}
\end{equation}
Their population MSE is $\sigma^2(1+\rho)/(3+\rho)$; equal weights give
$\sigma^2(3+2\rho)/9$. These are mathematical reference values, not measured
Risk RF outcomes. With the actual regularizers and delayed estimated matrix,
improvement must be measured rather than inferred from this reference.

Prespecified primary outcomes are prequential fused MSE and the paired
full-matrix-minus-diagonal difference within each $\rho$ condition. Secondary
outcomes are MAE, weights, cap activity, matrix-estimation error, and the fraction
of windows in which the methods choose different weights. A stationary segment
and a segment with a changed error relationship will test estimation delay and
inertia. A positive claim requires a useful held-out error reduction under
correlated or cancelling errors; merely updating off-diagonal entries or changing
weights is insufficient. This forecast test alone does not establish service value.

\subsubsection{Comparing admission at matched false-admission rates}
\label{sec:validation-gates}

The gate study will compare raw-score disagreement, common-target disagreement,
and a gate without disagreement on the same replay observations. Forecasts,
weights, missing-source masks, resource feasibility, worker eligibility, and
event definitions will be shared when isolating the gate. Thresholds will be
selected only on a separate historical validation stream to target the
prespecified nominal admission rates $1\%$ and $5\%$. No test event or test outcome
will be used to choose a threshold. When a discrete gate cannot attain the target
rate, its realized rate and the attainable operating points will be reported.

The held-out stream will include randomized event phases relative to window
boundaries, gradual and abrupt changes, operating-only changes, stale feedback,
source conflict, partial outages, and simultaneous outages. Admission coverage
will be reported at the matched nominal rate, together with conditional delay,
missed-event counts, training eligibility, and nominal admissions. Raw alarms,
admitted candidates, launched workers, and completed deployments will remain
separate counters. Complete evidence loss must cause explicit abstention.
An $8/10$ to $10/10$ comparison at one shared numerical threshold cannot, by
itself, exclude the possibility that the gain comes from a more permissive gate.

\subsubsection{Testing recoverability and net executed service benefit}
\label{sec:validation-service}

\paragraph{Recoverability before trigger comparison.}
For each backbone and drift mechanism, the first experiment will compare a
frozen policy with an ample-update policy and an oracle-timing update schedule,
while keeping the initial parameters, inference allocation protocol, observations,
data availability, optimizer, worker delays, and profile definitions fixed.
The oracle schedule may use the injected onset times as a declared privileged
reference; it is not an implementable online comparator. The ample-update
budget will be selected on calibration runs and specified before evaluation.
Learning curves will report performance before the change, immediately after
the change, and after each actual deployment. If updating cannot improve the
deployed policy under these controls, that condition does not establish a useful
retraining problem. If it can, later trigger comparisons must explain any failure
to realize the available recovery.

\paragraph{Allocation, timing, and expenditure controls.}
The no-retraining control will retain adaptive inference, and the fixed-profile
control will disable it. Their difference estimates the benefit of inference
allocation under the tested service model. RF, DA-RF, observed-loss, periodic,
and randomized-trigger arms will then use the same initial policies, training
profiles, update counts per worker, optimizer, eligibility rules, deployment
delays, and external arrivals and capacities. Both a common maximum budget and
actual costs will be reported. A count- and cost-matched timing experiment will
hold the number and types of completed updates fixed while varying launch times;
training compute, memory occupancy, and direct costs will be accounted for over
the entire active interval. Feasibility and unspent budgets will be reported.

To separate policy recovery from allocation-mediated effects, a complementary
controlled experiment will use a prespecified inference schedule with capacity
reserved for every declared training plan. The full closed-loop experiment will
retain the actual competition between training and inference. These are different
estimands and will not be pooled. If fusion variables enter the service-policy
input, all arms in an update-isolation experiment will receive the same observer
features, or that feature pathway will be included as a separate factorial
treatment. A change in backbone observations must not be silently attributed
to a change in retraining timing.

\paragraph{Primary endpoint and claim criterion.}
The primary service endpoint will be paired environment return relative to
adaptive no retraining, with its reward definition and training-cost accounting
fixed before evaluation. Completion, deadline loss, throughput, unsuccessful
requests, time-to-recovery, deployment count, and resource consumption will be
reported separately. A net-service-benefit claim for an extension requires a
positive paired return effect, uncertainty consistent with that claim, and no
violation of the declared resource or service constraints. Superiority over a
fixed inference profile, lower training cost alone, or a larger proxy value does
not meet that criterion. The classification study will likewise report paired
accuracy differences alongside NLL and ECE, use a declared calibration protocol,
and retain the distinction between prediction before label arrival and updates
after label arrival.

\subsubsection{Component activation and mechanism audit}
\label{sec:validation-mechanisms}

Each component will receive a perturbation designed to activate its mechanism:
reduced valid support for $q_N$, delayed timestamps for $q_A$, temporally unstable
telemetry for $q_V$, changed cross-source error relationships for $R_k$, rapid
quality changes for inertia, a dominant source for the cap, and partial or total
source loss for coverage. Paired ablations will preserve all other settings.
For each completed window and decision slot, the saved audit will include raw
scores, source forecasts, availability, support, age, bootstrap variability,
quality reference, residual matrix, weights, cap activity, disagreement,
coverage, gate status and rejection reason, feasible candidates, selected profiles,
worker launch, actual gradient steps, deployment time, and executed outcomes.

The audit will distinguish four stages: a perturbation activates a component;
the component changes weights or uncertainty; the change alters an admitted or
executed action; and the action improves a prespecified outcome after its cost.
Activation alone proves none of the later stages. Weight changes without action
changes demonstrate a representation effect; action changes without outcome
improvement demonstrate a control effect without an established benefit.
An inactive ablation is reported as inconclusive in that condition rather than
as proof that the component is either necessary or unnecessary. No new online-adaptation or queue-service numerical result will be added until
these protocols have been executed and their saved predictions and deployment
logs have been verified. The known-moment verification reported separately does
not satisfy that requirement.



\section{Claim discipline and submission readiness}

The defensible theoretical claim is conditional: at a fixed feasible set, the
fusion optimizer is stable to declared quality, risk, and anchor perturbations;
matched-marginal redundant sources can produce a strict joint-risk advantage;
and that advantage improves an explicitly defined threshold decision.
Neither these results nor the small mechanism simulation proves general queue
recovery, better classification calibration, or universal superiority.

A claim of useful service adaptation requires the proposed DA-RF implementation
to recover a policy in a condition where a matched updating intervention is
shown to work, and to produce a positive net executed-service effect over
adaptive no retraining after all active training costs. The inherited hidden-
relation failure is retained until a newly executed and transparently versioned
treatment changes that result. Components whose removal changes no action or
outcome remain negative or inactive findings rather than manufactured
contributions.

Before submission, replace prospective protocols by measured results only
where execution artifacts support them, complete the real-data and matched
budget comparisons, identify source code and data versions, and reconcile all
figures and notation with Eq.~\eqref{eq:main-fusion}. Information Fusion welcomes
fusion architectures, algorithms, and applications; editorial suitability
still depends on demonstrated significance, soundness, and reproducibility.
No acceptance probability or journal-standard certification follows from this
revision alone.

